\documentclass[a4paper,11pt]{article}
\usepackage[colorlinks=true, linkcolor=blue]{hyperref}
\usepackage[utf8]{inputenc}
\usepackage{amsmath, amssymb}
\usepackage{graphicx}
\usepackage{listings}
\usepackage{xcolor}
\usepackage{geometry}
\geometry{margin=1in}
\graphicspath{{img/}}
\title{Cubic Spline Interpolation}
\author{Ori Gonen}
\date{}

\lstset{
    language=Julia,
    basicstyle=\ttfamily\small,
    keywordstyle=\color{blue},
    commentstyle=\color{gray},
    stringstyle=\color{green!60!black},
    breaklines=true,
    showstringspaces=false
}

\begin{document}

\maketitle

\section{Introduction}

The goal of this assignment is to implement cubic spline interpolation. Given a set of points $(x_i, y_i)$, we construct a smooth function $S(x)$ that interpolates the data and is twice continuously differentiable.
\section{Mathematical Background}

On each interval $[x_i, x_{i+1}]$, the spline is a cubic polynomial:
\[
S_i(x) = a_i + b_i (x - x_i) + c_i (x - x_i)^2 + d_i (x - x_i)^3
\]

Its derivatives are:
\[
S_i'(x) = b_i + 2c_i(x - x_i) + 3d_i(x - x_i)^2
\]
\[
S_i''(x) = 2c_i + 6d_i(x - x_i)
\]

\subsection*{Conditions}
\begin{itemize}
    \item Interpolation:
    \[
    S_i(x_i) = y_i, \quad S_i(x_{i+1}) = y_{i+1}
    \]

    \item First derivative continuity:
    \[
    S_i'(x_{i+1}) = S_{i+1}'(x_{i+1})
    \]

    \item Second derivative continuity:
    \[
    S_i''(x_{i+1}) = S_{i+1}''(x_{i+1})
    \]

    \item Natural boundary conditions:
    \[
    S''(x_1) = 0, \quad S''(x_n) = 0
    \]
\end{itemize}

We get a linear system with $4(n-1)$ variables (the variables are the coefficients) and equations.

\subsection{Derivation of the Linear System}

We construct the spline on each interval $[x_i, x_{i+1}]$ as:
\begin{align*}
S_i(x) &= a_i + b_i(x - x_i) + c_i(x - x_i)^2 + d_i(x - x_i)^3 \\
S_i'(x) &= b_i + 2c_i(x - x_i) + 3d_i(x - x_i)^2 \\
S_i''(x) &= 2c_i + 6d_i(x - x_i)
\end{align*}

\paragraph{Natural boundary conditions}

At the left endpoint:
\begin{align*}
S''(x_1) &= 0 \\
2c_1 + 6d_1(x_1 - x_1) &= 0 \\
c_1 &= 0
\end{align*}

At the right endpoint:
\begin{align*}
S''(x_n) &= 0 \\
2c_{n-1} + 6d_{n-1}(x_n - x_{n-1}) &= 0 \\
c_{n-1} + 3d_{n-1}(x_n - x_{n-1}) &= 0
\end{align*}

\paragraph{Interpolation conditions}

For $i = 1, \dots, n-1$:
\begin{align*}
S_i(x_i) &= y_i \quad \Rightarrow \quad a_i = y_i \\
S_i(x_{i+1}) &= y_{i+1}
\end{align*}

Let $h_i = x_{i+1} - x_i$. Then:
\begin{align*}
a_i + b_i h_i + c_i h_i^2 + d_i h_i^3 = y_{i+1}
\end{align*}

\paragraph{First derivative continuity}

For $i = 1, \dots, n-2$:
\begin{align*}
S_i'(x_{i+1}) &= S_{i+1}'(x_{i+1}) \\
b_i + 2c_i h_i + 3d_i h_i^2 &= b_{i+1} \\
b_i + 2c_i h_i + 3d_i h_i^2 - b_{i+1} &= 0
\end{align*}

\paragraph{Second derivative continuity}

For $i = 1, \dots, n-2$:
\begin{align*}
S_i''(x_{i+1}) &= S_{i+1}''(x_{i+1}) \\
2c_i + 6d_i h_i &= 2c_{i+1} \\
c_i + 3d_i h_i - c_{i+1} &= 0
\end{align*}

\paragraph{Total number of equations}

We obtain:
\begin{itemize}
    \item 2 boundary equations
    \item $n-1$ interpolation equations at left point 
    \item $n-1$ interpolation equations at right point
    \item $n-2$ first derivative continuity equations
    \item $n-2$ second derivative continuity equations
\end{itemize}

Thus the total number of equations is
\[
2 + (n-1) + (n-1) + (n-2) + (n-2) = 4(n-1)
\]

This is also the number of coefficients! So we can solve the system.
\section{Implementation}

\subsection{Data Structure}

We define a struct:
\begin{lstlisting}
struct Zlepek
    x::Vector{Float64}
    y::Vector{Float64}
    coefficients::Vector{Vector{Float64}}
end
\end{lstlisting}

Each interval has coefficients $(a_i, b_i, c_i, d_i)$.

\subsection{Building the System}
construct a linear system:
\[
A \mathbf{u} = \mathbf{b}
\]

where $\mathbf{u}$ is a vector of length $4(n-1)$ which has all the coefficients. \\ 
We then solve by Julia's backslash operator.
The solution vector is then reshaped into per-interval coefficients. (vector of $4(n-1)$ floats $\longrightarrow$ vector of $n-1$ vectors of 4 coefficients)

\section{Evaluation}

To evaluate the spline at a point $x$, we:
\begin{itemize}
    \item Find the interval index $i$ where $x\in [x_i,x_{i+1}]$.
    \item Evaluate with the corresponding coefficients
\end{itemize}



\section{Testing}

We verify correctness using multiple functions:
\begin{itemize}
    \item Linear function
    \item Constant
    \item sine
    \item sine with uneven spacing
\end{itemize}



\section{Results}

The implementation successfully:
\begin{itemize}
    \item Interpolates all given points
    \item Produces smooth curves
\end{itemize}
Different example functions (sin, random point, Gaussian) are presented in demo file, and can be seen in Figure \ref{fig:demo}.
\begin{figure}[h]
    \centering
    \includegraphics[width=0.8\textwidth]{img/demo.png}
    \caption{Demonstrations with different functions}
    \label{fig:demo}
\end{figure}
\\ \\
Downfall: undersampling (i.e., points are too far apart for "high frequency" functions) results in interpolation which is not representative of the original function. As we can see in Figure \ref{fig:undersamp}, not sampling sufficiently frequent will result in unrepresentative interpolations. In the top row, we sample frequent enough to not tell the difference. however, on the bottom we see problems: in the left we see we always sample on 0s of the function, so the interpolation is constant 0. on the right, we sample so slow that it looks like we are sampling a lower frequency \textsf{sin} function.
\begin{figure}[h]
    \centering
    \includegraphics[width=0.8\textwidth]{img/undersampling_demo.png}
    \caption{Cubic spline interpolation for Sin function with different sampling frequencies}
    \label{fig:undersamp}
\end{figure}




\end{document}